% Fragmento público generado desde propietarios íntegros.
% Residencia: 52.7.1.
% Fuente: ../incoming/radion_monografia_20260827/manuscrito/sections/08_complejo_electromagnetismo.tex:1-58
\noindent Los 3 regímenes del TRIT admiten realizaciones compleja, electromagnética y mecánica con tipos diferenciados.\par\medskip

\subsubsection{Operador cuadrático de raíz interna de \(-1\)}

\paragraph{Representación exponencial del régimen elíptico \(J^2=-I\)}

Con \(J^2=-I\),
\begin{equation}
T_{\mathrm{ell}}(x,y)=e^{-xI-yJ}
=e^{-x}(\cos y\,I-\sin y\,J).
\label{eq:elliptic-exp}
\end{equation}
Bajo la identificación generada por \(J\),
\[
T_{\mathrm{ell}}(x,y)\longleftrightarrow e^{-x-iy}.
\]
El escalar \(x\) controla dilatación; \(y\) controla giro orientado. La forma
\(a+bi\) se reconoce como elasticidad más rotación, no como dos cantidades
pegadas arbitrariamente.

\paragraph{Representación exponencial del régimen parabólico \(J^2=0\)}

Con \(N^2=0\), la exponencial se trunca:
\begin{equation}
T_{\mathrm{par}}(x,y)=e^{-x}(I-yN).
\label{eq:parabolic-exp}
\end{equation}
El régimen describe el umbral donde no hay ni giro periódico ni apertura
exponencial. En la lectura temporal activa es el presente euclídeo, pero
conserva memoria de las ramas que lo limitan.

\paragraph{Representación exponencial del régimen hiperbólico \(J^2=+I\)}

Con \(R^2=I\) y \(P_\pm=(I\pm R)/2\),
\begin{equation}
T_{\mathrm{hyp}}(x,y)=e^{-xI-yR}
=q_+P_++q_-P_-,
\qquad q_\pm=e^{-x\mp y}.
\label{eq:hyperbolic-exp}
\end{equation}
Los invariantes elementales son
\begin{equation}
q_+q_-=e^{-2x},
\qquad
\frac{q_-}{q_+}=e^{2y}.
\label{eq:q-product-ratio}
\end{equation}
El producto conserva el modo común; el cociente conserva la separación
orientada.

Aplicado al par angular,
\[
x=\frac{\pi A}{180},\qquad y=\frac{\pi C^*}{180}.
\]
Por ello, \(A\) es la coordenada común de la publicación y \(C^*\) su
contraste orientado. Sólo después de aplicar un lector físico se reconocen
como polos eléctricos y magnéticos.

